\section{总结与局限}
\label{sec:conclusion}
\subsection{项目工作总结}
本项目在统一 Breakout 环境中实现了 DQN 与 PPO 的图像控制、训练、开发选模和实时展示。两者共用动作、奖励、回合终止及观测接口；DQN 以回放和目标网络学习动作价值，PPO 以 GAE 与裁剪目标优化动作分布。日志、检查点和逐局分数使实验可以复核。

DQN 低学习率续训得到较好的中间模型，继续训练和后续恢复却未进一步提高得分。PPO 的额外价值拟合降低了训练误差，但未改善开发所选策略；提高 GAE $\lambda$ 的额外收益也未得到区间支持。新增 KL 对照虽实际提高了策略学习率和整轮 KL，三次续训的测试收益仍有正有负，总体区间跨零。两条链路的共同经验是保留负结果、依据开发集选模，并用独立测试检查策略表现。

\subsection{局限与后续验证}
原 DQN 的配对复测仍属于开发数据。本次在冻结模型后，使用留出的共同开局完成同协议的 50 局测试及随机基线，已在第\ref{sec:comparison}节集中比较；未依据这些测试更新权重或更换主比较候选。结果仅支持固定模型在当前规则下的比较，不能证明稳定通关或跨环境泛化。

本实验并非同预算、同初始化的算法排名：DQN 与 PPO 使用不同的多阶段训练历史和超参数搜索，既不能用单次测试差推断普遍优劣，也不能从当前耗时推断公平的样本效率。更强结论需要统一从头训练预算、多个独立训练种子和未用于调参的新测试集。当前最可靠的结论是两个已训练系统在同一任务上的可复核表现，以及训练、选模和稳定性取舍的实测差异。

后续需对入围设置进行多个独立训练种子的重复实验，并保留未用于调参的测试集。除平均分外，还应报告标准差、中位数、低分位数、低分比例、通关率和时间截断比例。对同一模型增加开局种子，不能替代独立重复训练。

对于学习率、回报跨度和回放容量的作用，可以从同一模型起点出发，固定其余设置并分别改变一个因素。这样能够进一步区分模型起点、额外交互预算和参数设置对结果的影响。当前恢复实验的负结果应保留为已有观察，而不是在缺少新数据时替换为预期提升。

\FloatBarrier
